\section{Cross-paper comparison: do not collapse the design choices}\label{sec:comparison}
\subsection{Where forward states and actor updates come from}
\begin{longtable}{@{}L{0.16\linewidth}L{0.25\linewidth}L{0.26\linewidth}L{0.25\linewidth}@{}}
\toprule Method & Forward data for control learning & Policy update route & Model's main role\\\midrule\endhead
PETS & Predicted physical-state particles & Action-sequence search & Uncertainty-aware planning\\
MBPO & Short model branches from real replay & SAC actor--critic & Synthetic transition data\\
Dreamer & Predicted latent trajectories & Differentiated imagined values & Latent inference and imagination\\
ContactNets & Mechanics-based contact transitions & No single policy-learning recipe & Structured contact identification\\
SHAC & Differentiable simulator windows & Pathwise simulator gradient & Known numerical dynamics\\
PWM & Predicted task-conditioned latents & Pathwise model gradient & Pretrained optimization interface\\
DMO & Simulator trajectories & Learned model Jacobians backward & Local sensitivity approximation\\
RWM-U & Predicted observation histories & PPO with uncertainty penalty & Offline neural simulator\\
SuperMimic & Predicted robot states from replay starts & Pathwise model gradient & Online shared motion dynamics\\
$\nabla$Mimic & Specialized simulator windows & Clipped action gradients and target regression & Simulator adjoint, no learned dynamics\\
WM-LOCO & Simulator interaction; predictive recurrent memory & PPO actor--critic & Terrain-aware predictive context\\
\bottomrule
\end{longtable}
The entries separate algorithmic routes; they are not a performance ranking.
A method can share an RSSM with another while using a different actor estimator.
It can also share a short-horizon value-tail objective while obtaining derivatives
from a different source.

\subsection{The evidence relevant to geometry adaptation}
\begin{longtable}{@{}L{0.16\linewidth}L{0.40\linewidth}L{0.36\linewidth}@{}}
\toprule Method & Relevant evidence in the reviewed work & Transfer question still requiring a test\\\midrule\endhead
PETS & Uncertainty-aware planning on control benchmarks & Does disagreement detect unseen robot--terrain interactions?\\
MBPO & Short model branches for online control learning & How should trusted horizon change at new contacts?\\
Dreamer & Visual control from latent imagination & Does a source representation preserve new support geometry?\\
ContactNets & Contact identification from real transitions & Does structure lower humanoid target-data cost?\\
SHAC & Differentiable locomotion/control benchmarks & What remains when learned-model bias is removed?\\
PWM & Regularized model-based policy extraction & Does source-model optimization structure survive a geometry shift?\\
DMO & Simulation benchmarks and Go2 locomotion deployment & Is target failure caused by forward drift or local sensitivities?\\
RWM-U & Offline robotic policy learning and deployment & When does caution need to become new-data collection?\\
SuperMimic & Broad motion tracking and G1 deployment & How much dynamics knowledge transfers beyond source contacts?\\
$\nabla$Mimic & Simulator-gradient G1 tracking & Do finite contact-changing updates remain useful on new terrain?\\
WM-LOCO & Visual humanoid traversal of constrained footholds & Does motion pretraining improve data-to-success or model gradients?\\
\bottomrule
\end{longtable}
These questions are review interpretations, not claims that every paper attempted
and failed the same experiment. The cited tasks, observations, rewards, hardware,
seed counts, and clocks differ. No common sample-efficiency leaderboard can be
constructed from their headline scores alone.

\subsection{Three axes that should remain independent}
\begin{enumerate}
\item \textbf{Information}: physical state, proprioceptive history, camera/depth,
known geometry, and hidden material context.
\item \textbf{Representation}: state residuals, recurrent latents, ensemble
predictions, or mechanics/contact structure.
\item \textbf{Optimization}: sampling-based planning, model-generated data,
pathwise imagined return, decoupled gradients, or model-conditioned PPO.
\end{enumerate}
A new method changes a scientifically identifiable component only if the other
axes are controlled or their effects are measured. A large recurrent vision
policy with new foothold rewards is not a clean comparison to a small
proprioceptive motion tracker on a different objective.

\section{The precise research gap and what would count as an answer}\label{sec:gap}
\subsection{Define the source and target distributions}
Let $\mathcal D_S$ be source transitions collected during motion tracking on
flat ground, and $\phi_0$ the resulting learned model. A target environment is
sampled from a declared distribution $E\sim\mathcal E_T$, such as stairs with
specified ranges of riser height, step depth, orientation, and approach pose.
Material parameters $\xi$ are initially held fixed to isolate geometry.
Motion pretraining here means fitting executed state--action--successor
transitions collected during tracking. Reference poses alone do not identify
the physical consequences of intervening with different motor commands.

After $N$ additional executed target transitions,
\begin{equation}
\phi_N=\operatorname{Adapt}(\phi_0,\mathcal D_S,\mathcal D_T^{(N)}),\qquad
\theta_N=\operatorname{Improve}(\theta_0,\phi_N,\mathcal D_T^{(N)}).
\end{equation}
This notation does not assume that model and policy adaptation happen in two
separate phases. They can alternate. It does require that the counted target
interaction includes all data collection induced by either process.

For a fixed held-out target test distribution, define
\begin{equation}
\mathcal S(N)=\Pr[\text{successful task completion under }\pi_{\theta_N}],\qquad
N_\tau=\inf\{N:\operatorname{LCB}(\mathcal S(N))\ge\tau\}.
\end{equation}
$\operatorname{LCB}$ is a declared confidence lower bound, and $\tau$ is a
preselected success target. If the budget is exhausted first, $N_\tau$ is
unreached/censored, not the final budget labeled as a success.

\subsection{Why the literature does not give a universal interaction count}
Changing the target geometry's range changes what has to be learned. Supplying
an exact map removes a perception problem that a depth-only agent must solve.
Initializing a competent policy can make data collection much easier. Reward
shaping can turn a sparse task into a dense one. Different simulator integration
rates also make a raw ``step'' count ambiguous.

The defensible answer is therefore a curve $\mathcal S(N)$ plus its experimental
conditions, not an unsupported prediction such as ``a few thousand samples
should be enough.'' An interaction count without source coverage, target
distribution, and a controller specification is not a transferable result.

Report both adaptation and total data cost:
\begin{equation}
N_{\rm total}=N_{\rm pre}+\sum_{j=1}^{J}N_{\rm new}^{(j)},\qquad
N_{\rm amortized}=N_{\rm pre}/J+\frac1J\sum_jN_{\rm new}^{(j)}.
\end{equation}
Amortization across $J$ target tasks is a legitimate motivation, but it should
not hide pretraining cost when comparing with a from-scratch learner.

\subsection{Separate geometry, material, and task transfer}
\begin{tabularx}{\linewidth}{@{}lXX@{}}
\toprule Shift & First controlled example & New information potentially needed\\
\midrule
Geometry & One step, then varying stairs & Surface locations, normals, feasible support\\
Material & Same stair shape with varied friction & History-dependent material inference or supplied $\xi$\\
Interaction & Hands on a rail or moving objects & New contact pairs and object state/constraints\\
Task & Tracking versus climb-to-goal & Goals, reward, termination, and value adaptation\\
\bottomrule
\end{tabularx}
This ordering is a proposed research progression. It is not a claim that a
height map can represent every later interaction. A two-and-a-half-dimensional
terrain map can be useful for stairs while failing to represent overhangs or a
handrail's full geometry.

\subsection{What would be a substantive contribution?}
Possible findings include a reproducible reduction in target interaction for
fixed success, an explanation of which source motions provide reusable contact
coverage, or a structure/regularizer that improves actual update utility without
simply increasing model size or privileged information.

Adding an ensemble, residual adapter, or terrain encoder is not by itself a
new scientific contribution. Those ingredients have precedents. The contribution
would be the measured transfer behavior and a convincing explanation of why it
occurs under new contact geometry. This focused selection also does not prove
novelty against every adjacent paper; a final priority search remains necessary
before making a publication claim.

\section{Model structures worth comparing}\label{sec:structures}
\subsection{Conditioned state model: the minimal baseline}
The smallest informative extension is
\begin{equation}
\hat x_{t+1}=\operatorname{Integrate}\big(x_t,
\delta_\phi(x_t,a_t,e_t,\xi)\big).
\end{equation}
The terrain query is evaluated consistently along the predicted pose. Start with
known geometry and state measurements, so prediction errors cannot be blamed
on a camera or state estimator. This baseline tests whether initialization and
conditioning already deliver most of the transfer benefit.

A zero-initialized new conditioning pathway can preserve the old model's initial
behavior while exposing terrain information. Preservation of outputs is not
evidence that those outputs are correct on stairs; it is a controlled starting
condition for adaptation.

\subsection{A residual adapter in the correct state coordinates}
An alternative keeps a pretrained prediction and learns a correction:
\begin{align}
\delta_N(x,a,e)&=\delta_{\phi_0}(x,a)+\delta_{\phi_A}^{\rm corr}(x,a,e),\\
\hat x'&=\operatorname{Integrate}(x,\delta_N).
\end{align}
The adapter has its own parameters $\phi_A$; they are trainable dynamics
parameters, not a critic. Corrections should be applied in a valid residual or
tangent representation rather than added directly to quaternion coordinates.

The base already contains flat-ground contact responses. The adapter may need
to \emph{cancel} an incorrect old collision prediction as well as add a new one.
It is not automatically a clean decomposition into universal robot physics plus
environment forces. Compare frozen-base adaptation, partial unfreezing, and
full fine-tuning with explicit trainable-parameter and compute budgets.

\subsection{Known mechanics with learned contact effects}
Use known mass, actuation, and robot geometry with a terrain query, then learn
contact quantities or unknown corrections:
\begin{equation}
v'=v+M^{-1}\left[\Delta t\,f_{\rm free}(x,\mathcal U(a))
+J(q,E)^Tj_\phi(x,a,e,\xi)\right].
\end{equation}
The function producing $j_\phi$ could be an implicit resolver, a constrained
network, or a learned residual around a known solver. These are different
models and must not be represented by the same vague label ``physics informed.''

Potential advantages are reuse of known geometry and fewer physically impossible
predictions. Potential disadvantages include wrong free dynamics, contact-set
discovery, computational expense, and unstable derivatives of the resolver.
Known structure can move the difficulty rather than eliminate it.

\subsection{Recurrent or latent variants}
If the target is partially observed, a recurrent state can represent hidden
velocities, contact history, or material context. This adds a belief-estimation
problem to dynamics adaptation. A strong first experiment should compare it
with a fully observed geometry-conditioned model, not change observations and
model class simultaneously.

Stochastic latents and ensembles can describe ambiguity, but they do not supply
missing geometric observations. A unimodal conditional mean can be physically
implausible when unobserved contact outcomes are multimodal. A more expressive
distribution may represent those alternatives without identifying which one is
actually present; planning and data collection must then use that uncertainty.

\section{What makes a gradient useful?}\label{sec:gradientutility}
\subsection{Forward error and sensitivity error have different recurrences}
On a region where a fixed-policy transition has Lipschitz constant $L$, a
one-step forward error bound $\varepsilon_F$, with
$\varepsilon_t^x=\norm{\hat x_t-x_t}$, gives
\begin{equation}
\varepsilon_{t+1}^x\le L\varepsilon_t^x+\varepsilon_F.
\end{equation}
This is a local smooth-region argument, not a contact-mode-independent theorem.
For policy sensitivity $S_t=D_\theta x_t$, write
$A_t=F_x+F_aA_x$ and $B_t=F_aA_\theta$. Along aligned evaluation points,
\begin{align}
S_{t+1}&=A_tS_t+B_t,\\
\norm{\hat S_{t+1}-S_{t+1}}&\le\norm{\hat A_t}\norm{\hat S_t-S_t}
+\norm{\hat A_t-A_t}\norm{S_t}+\norm{\hat B_t-B_t}.
\end{align}
The second line is a direct norm inequality derived for this review. If the
models visit different points, those distribution differences also enter the
Jacobian mismatch. A state-fit bound alone does not bound the last two terms.

\subsection{Local derivatives can miss contact-changing moves}
Consider the synthetic one-dimensional unilateral template
$x'=\max(0,a)$ with target reward $-(x'-x^*)^2$. For $a<0$, the local derivative
is zero even though a large enough positive action change can improve the target.
A softplus surrogate supplies a nonzero direction, but changes forward outcomes
and can be optimistic. This is a toy constraint map, not a G1 stair simulation.

\begin{figure}[htbp]\centering
\includegraphics[width=\linewidth]{figures/contact_gradient.pdf}
\caption{Synthetic review illustration. A useful finite change can cross a mode
boundary even when a local derivative is zero. Smoothing can create a direction
while introducing bias. Neither curve is an empirical contact benchmark.}
\end{figure}

Not every contact event produces an unbounded derivative. Some yield bounded
but discontinuous derivatives, others a zero local sensitivity, and some
geometric singularities produce extreme values. A method designed only to
clip explosions does not address all three cases.

\subsection{Finite-update utility should be measured directly}
Let $d$ be a proposed policy update direction, normalized to a matched physical
action change on a fixed probe distribution. Measure
\begin{align}
\Delta J_{\rm sim}(\epsilon,d)&=J_{\rm sim}(\theta+\epsilon d)-J_{\rm sim}(\theta),\\
D_{\epsilon}(d)&=\frac{J_{\rm sim}(\theta+\epsilon d)-J_{\rm sim}(\theta-\epsilon d)}{2\epsilon}.
\end{align}
The second quantity is a finite-change score. At contact boundaries it should
not be advertised as an exact derivative. Sweep $\epsilon$ over meaningful
action scales and use matched starting states, noise realizations, and terrain
instances to reduce comparison variance.

Useful summaries include the fraction of predicted improvements that improve
actual return, the gain at a fixed action-change budget, and agreement between
predicted and observed improvement. Report these separately for contact events,
free motion, and contact-mode changes. A cosine against a tiny finite difference
alone is insufficient.

\subsection{A critic is part of the optimization model}
Poor value extrapolation can make an accurate dynamics model look bad, and a
good critic can compensate for a short horizon. Freeze or standardize the critic
where feasible in diagnostic comparisons, then test full joint learning.
Separate failures caused by transition prediction, transition sensitivities,
reward representation, continuation, and tail value.

For example, compare an identical rollout with an analytic task reward and a
learned reward head. If their gradients differ near a support transition, the
problem may be reward prediction rather than the transition model. The same
logic applies to differentiable termination surrogates.

\section{A staged experimental program}\label{sec:experiments}
\subsection{Stage 0: establish an auditable source baseline}
Verify joint ordering, action units, state frames, geometry, reset behavior, and
terminal successors. Evaluate the source tracker over full clips, not only short
prefixes. Save exact model and policy checkpoints. Record source transitions,
training compute, and evaluation settings before claiming transfer benefits.

Our local motion model can be a source checkpoint for this study, but its
existing five-second training evaluations and 20-second video prefixes do not
establish a matched reproduction of published full-motion scores. That limitation
should remain visible in any subsequent experiment report.

\subsection{Stage 1: isolate prediction transfer on one known step}
Use exact simulator geometry, fixed material parameters, and a single low step.
Collect a declared set of approach, swing, touchdown, toe-strike, and recovery
transitions. Keep a held-out set of contact situations, with entire trajectories
separated from training rather than randomly splitting adjacent frames.

Compare random initialization, full source-model fine-tuning, and a source model
with a new terrain adapter. Train with the same target data at several budgets.
Evaluate multistep state error, contact timing, surface penetration, support
location, and sensitivity probes. Match normalization and input information.

\subsection{Stage 2: separate world-model reuse from policy reuse}
\begin{center}
\begin{tabular}{@{}lll@{}}
\toprule Condition & Dynamics initialization & Policy initialization\\\midrule
A & Random & Random\\
B & Motion pretrained & Random\\
C & Random & Motion pretrained\\
D & Motion pretrained & Motion pretrained\\
\bottomrule
\end{tabular}
\end{center}
Use the same target controller input, reward, geometry distribution, and reset
curriculum. This factorial design distinguishes a better dynamics prior from
a policy that simply collects better data. A frozen source model is an additional
diagnostic, not necessarily a competitive final method.

If an old policy is extended with new terrain inputs, specify how those inputs
are initialized and which weights are trainable. Otherwise, an architecture
change and a pretraining change are entangled.

\subsection{Stage 3: cross structure with optimization route}
For the most promising models, compare fully imagined BPTT, simulator-forward
model-backward updates, and a model-generated-data or predictive-feature baseline.
Use matched policy/value capacity when possible. A conventional terrain-aware
PPO controller provides a model-free reference under the same task.

Compare these routes within each model structure, rather than replacing both
the controller and model simultaneously. A forward/backward comparison can
identify rollout drift; a BPTT/PPO comparison can identify sensitivity dependence;
a plain/structured comparison can identify representation effects.

\subsection{Stage 4: introduce genuine distribution shifts}
Expand from one step to held-out stair heights, step depths, orientations, and
descent. Hold out combinations as well as individual parameter values. Then add
depth-based terrain inference, occlusion, and material variation one at a time.
Finally consider new contact bodies or moving objects.

Always distinguish interpolation inside a randomized training range from
extrapolation outside that range. A new terrain instance from the same generator
is not the same transfer test as a new collision topology.

\subsection{Measurement and accounting protocol}
\begin{enumerate}
\item Count transitions across all worlds. Include target warm-up, failures,
resets that generate learning data, and simulator-forward gradient windows.
Keep physics substeps and control transitions separate.
Record evaluation interactions separately, including their cost when evaluation
is used for selection or when reporting total robot interaction.
\item Report data-to-success curves with multiple training seeds and separate
terrain instances. Use a fixed evaluation protocol and confidence intervals.
\item Report wall-clock time, model-update work, policy-update work, and
pretraining cost alongside target interaction.
\item Measure old-skill retention and new-task success, including failed runs.
Do not select a favorable checkpoint using the final test set without reporting it.
\item Report naturally distributed error and contact-balanced error. Average
MSE can hide rare interactions that determine task success.
\item Calibrate uncertainty and update-utility diagnostics on validation data;
reserve separate tests for final claims.
\end{enumerate}

\subsection{What outcomes would change our beliefs?}
If pretraining helps only when the policy is also pretrained, the evidence for
reusable \emph{dynamics} knowledge is weaker than a combined-system gain suggests.
If structure reduces physical violations but not data-to-success, it may be
improving a secondary prediction metric rather than the controller's learning signal.

If decoupled updates repair failures of imagined BPTT, prioritize forward model
coverage and rollout trust. If neither works but model-conditioned PPO does,
prioritize derivative quality or contact-mode exploration. If all methods fail
with the same reference objective, first inspect the task specification.

If a plain conditioned model fine-tunes as efficiently as a structured model,
the additional structure has not justified its cost under that target distribution.
Conversely, a benefit that persists across held-out contact geometries and matched
data budgets is much stronger evidence than a single successful staircase video.
